\chapter{SCADA and HMI Standards}
\section{Ignition Architecture}
Ignition is the preferred SCADA/HMI platform for new projects. The project team should choose the Gateway, database, client, and redundancy arrangement based on system size, availability needs, performance, and recovery time. Operator workstations should normally run Perspective or Vision clients rather than host the production Gateway or database.

Record the tag providers, projects, databases, historian, alarm pipelines, identity provider, certificates, modules, backups, and important external dependencies. Use a development or test environment when practical instead of making untested changes directly in production.

\section{Project and Tag Design}
Use shared projects, inheritance, UDTs, templates, and common scripts to avoid duplicating the same work across projects. Tag names and folders should follow the PLC equipment model and use consistent units, ranges, descriptions, quality handling, alarms, and history settings.

Keep scripts readable and reusable. Log errors that will help with troubleshooting, but avoid repetitive logging that hides useful information. Where a tag, expression, binding, or shared function can do the job cleanly, prefer it over copying scripts into many components.

\section{HMI Philosophy}
Screens should make the current process condition easy to understand. Use color mainly to show abnormal conditions or items requiring attention. Do not rely on color alone; use text, symbols, and state indication as well. Navigation, faceplates, trends, units, decimal precision, manual controls, and alarm presentation should work the same way throughout the project.

Use a clear display hierarchy, normally including overview, area, unit, equipment-detail, and diagnostic screens. Operators should be able to see the active mode, command ownership, equipment state, alarms, interlocks, and bypasses before issuing a command. Use confirmation prompts for consequential actions without adding unnecessary prompts to routine work.

\section{Perspective Design}
Perspective views should be tested at the screen sizes and aspect ratios expected in service. Use responsive containers where resizing is required. Docked views, popups, embedded views, parameters, and message handlers should have clear interfaces and should not depend on fragile component paths.

Startup and navigation scripts should be safe to run more than once, finish quickly, and handle missing tags or bad-quality data without breaking the view. Reusable views should not assume that every optional parameter is present.

\section{Alarms}
Configure Ignition alarms from the approved alarm list. Include the identifier, priority, message, cause, consequence, expected response, deadband, delay, shelving permissions, and any state-based suppression. Acknowledging an alarm does not clear the process condition that caused it.

Test alarm journals and notification pipelines, including loss and recovery of the database or notification service. Review nuisance, chattering, standing, and frequently occurring alarms after startup.

\section{Historian and Data}
Historize data that has a defined operating, quality, maintenance, regulatory, or analysis purpose. Choose the sampling mode, deadband, timestamp source, and retention period to match that purpose. Document store-and-forward behavior and how bad or uncertain quality is handled.

Database connections should use managed credentials and only the permissions they need. Use encrypted connections when supported by the approved architecture.

\section{Security and Audit}
Use the approved identity provider when practical and assign access by role. Typical roles include operator, maintenance, engineering, supervisor, and administrator. Audit important actions such as setpoint changes, alarm shelving, recipe changes, and security administration.

OPC UA connections should use trusted certificates and supported security policies. Remove default accounts and unused connections before turnover.

\section{Backup, Deployment, and Recovery}
Create a Gateway backup before and after significant production changes and after commissioning. Record the Ignition version, installed modules, licenses, database dependencies, certificates, identity-provider settings, DNS/NTP needs, and restore steps.

For critical systems, demonstrate that the backup can be restored and that the project can be rolled back if a deployment fails. Keep the accepted backup with the project turnover package.
